\section{Introduction}\label{sec:introduction}

This article is a rigorous computational case study of the explicit
ten-state weighted synchronisation kernel displayed in
\Cref{def:sync-kernel} and Figure~\ref{fig:sync-kernel}.  The kernel has
uniform outdegree \(3\), a strongly connected aperiodic underlying
digraph, and a binary output labelling.  Its ten states split into five
dual pairs, and the duality exchanges the output-zero and output-one
edges.  These features make it a small but nontrivial test case for the
interaction between labelled finite-state dynamics, self-dual determinant
identities, plane-curve geometry, finite-field Galois certificates, and
cyclotomic twists.

The scope is intentionally concrete.  We do not classify self-dual
kernels, and we do not claim that the ten-state example represents a
generic family.  Instead, every assertion below is made for the displayed
kernel or for an explicitly stated auxiliary criterion that is then
verified for this kernel.  The point of the paper is that this one
kernel admits a complete audit trail: the transition table gives the
matrix, the matrix gives the determinant, the determinant gives the
completed curve and the cyclotomic factors, and the Sage runner
\[
  \texttt{certificates/verify\_certificates.py}
\]
reproduces the exact symbolic certificates cited in the proofs.

Let \(B(u)=B_0+uB_1\) be the weighted adjacency matrix of this kernel and
write
\[
  \Delta(z,u):=\det(I-zB(u)).
\]
The self-duality permutation satisfies
\[
  \Pi^{-1}B_0\Pi=B_1,\qquad \Pi^{-1}B_1\Pi=B_0,
\]
and therefore
\[
  \Delta(z,u)=\Delta(uz,u^{-1}).
\]
After the change of variables
\[
  r=e^{i\theta},\qquad u=r^2,\qquad
  s=r+r^{-1}=2\cos\theta,\qquad w=zr,
\]
the determinant becomes the completed determinant
\(\widehat\Delta(w,s)\in\ZZ[s][w]\).  The completion compresses the
unit-circle twists \(u=e^{2i\theta}\) to the single real parameter
\(s=2\cos\theta\), and in this example the resulting polynomial has
degree \(6\) in \(w\) although \(B(u)\) is a \(10\times10\) matrix.
This low-degree completion is the phenomenon isolated in the paper: the
same polynomial controls the quotient curve, the generic Galois group,
the primitive real cyclotomic factors, and the large-\(m\) twisted
Perron-root asymptotics.

The same weighted matrix also controls the natural cyclic skew-product
lifts
\[
  \widetilde B_q:=B_0\otimes I_q+B_1\otimes P_q,
\]
where \(P_q\) is the cyclic permutation matrix on \(\ZZ/q\ZZ\). Their
Fourier blocks are \(B(\omega_q^j)\), so the unit-circle twists are the
natural character factors of the labelled cyclic lifts. The cyclotomic
root products considered below are precisely the primitive real factors
extracted from this Fourier decomposition after the completed
coordinate change.
Throughout, such a factor is the conjugate product itself, or
equivalently the ordered resultant with the real cyclotomic polynomial
in the first argument. The opposite resultant order differs by the
usual Sylvester sign and is never used as the definition of \(R_m\).


\paragraph{Terminology.}
We call \((w,s)\mapsto \widehat\Delta(w,s)\) the completed determinant in analogy with completed \(L\)-functions: it is simply the polynomial obtained from \(\Delta(w/r,r^2)\) after rewriting in the self-dual invariants \(w\) and \(s=r+r^{-1}\). Likewise \(\Phi_m^+\) denotes the real cyclotomic polynomial of \(2\cos(\pi/m)\); concretely
\[
  \Phi_m^+(s)
  =\prod_{a\in(\ZZ/2m\ZZ)^\times/\{\pm1\}}
    \bigl(s-2\cos(\pi a/m)\bigr),
\]
so \(\Phi_m^+\) is the standard real Chebyshev factor attached to the primitive \((2m)\)-th roots of unity.



The case study has two strands. The first is geometric: the affine curve
\(\widehat\Delta(w,s)=0\) admits the sign involution
\((w,s)\mapsto(-w,-s)\), and its normalisation and quotient can be
described explicitly. The second is arithmetic: eliminating the
parameter \(s\) along cyclotomic slices produces one-variable root
products, equivalently ordered resultants. For this kernel, explicit
root-separation certificates show that the endpoint conjugates select
the minimal-modulus roots for large \(m\), giving a high-order
asymptotic expansion for the cyclotomic twists.

\paragraph{Computer-assisted verification.}
The exact determinant, discriminant, finite-field factorisations,
smoothness ideals, interval separation checks, and endpoint jet
coefficients were computed in SageMath~10.2 and cross-checked in
Mathematica~14.0 where indicated.  The appendix gives a
theorem-to-command map and the single Sage certificate runner
\(\texttt{certificates/verify\_certificates.py}\).  The runner prints
the Sage and Python versions, SHA-256 checksums for the paper sources
and the certificate script, and named outputs such as
\(\texttt{determinant\_difference: 0}\),
\(\texttt{self\_duality\_difference: 0}\),
\(\texttt{affine\_smoothness\_groebner: [1]}\),
\(\texttt{mod7\_s3\_gcd\_factor}\),
\(\texttt{mod19\_s3\_gcd\_factor}\),
\(\texttt{endpoint\_interval\_certificate: true}\), and
\(\texttt{finite\_m\_branch\_selection\_4\_to\_31: true}\).
Items \textup{(a)}--\textup{(f)} in
\Cref{sec:appendix-computation} display the exact outputs used in the
corresponding proofs; no fitted parameters are used.


\subsection{Main results}

The main results are kernel-specific certificates for the ten-state
matrix.  We first record the determinant, spectral radius, completion,
quotient, and generic Galois group.

All curves in the main results are considered over \(\overline{\QQ}\), and
\(\overline{X}\) denotes the normalisation of the projective closure of
the irreducible affine curve \(\widehat\Delta(w,s)=0\).

\begin{theorem}\label{thm:intro-kernel-geometry}
Let \(B(u)\) be the weighted synchronisation kernel of
\Cref{def:sync-kernel}. Then:
\begin{enumerate}[label=\textup{(\roman*)}]
  \item The determinant is
  \[
    \Delta(z,u)
    =1-(1+u)z-5uz^2+3u(1+u)z^3-u(u^2-3u+1)z^4
  \]
  \[
    \qquad
    +u(u^3-3u^2-3u+1)z^5+u^2(u^2+u+1)z^6.
  \]
  In particular,
  \[
    \Delta(z,1)=(z-1)(z+1)(3z-1)(z^3-z^2+z+1),
  \]
  Moreover \(B(1)\mathbf1=3\mathbf1\), hence
  \(\rho(B(1))=3\) by the row-sum bound, and
  \(C_B=243/272\).

  \item The completed determinant
  \[
    \widehat\Delta(w,s):=\Delta(w/r,r^2)
    \qquad (u=r^2,\ s=r+r^{-1})
  \]
  is given by
  \[
    \widehat\Delta(w,s)
    =1-sw-5w^2+3sw^3+(5-s^2)w^4+(s^3-6s)w^5+(s^2-1)w^6.
  \]
  It satisfies
  \[
    \widehat\Delta(w,-s)=\widehat\Delta(-w,s).
  \]

  \item The smooth projective curve \(\overline{X}\) has genus \(6\).
  Its quotient by the involution \((w,s)\mapsto(-w,-s)\) is a smooth
  plane quartic of genus \(3\), and the generic Galois group of
  \(\widehat\Delta(w,s)\) over \(\QQ(s)\) is \(S_6\).
\end{enumerate}
\end{theorem}

The second kernel-specific theorem concerns cyclotomic twists.

\begin{theorem}\label{thm:intro-kernel-cyclotomic}
For \(m\ge 2\), let \(s_m=2\cos(\pi/m)\), let \(\Phi_m^+\) be the
minimal polynomial of \(s_m\), let \(A_m\) be its real cyclotomic
conjugate set, and define
\[
  R_m(w):=\prod_{\alpha\in A_m}\widehat\Delta(w,\alpha).
\]
With the standard Sylvester convention this product is represented by
the ordered resultant
\[
  R_m(w)=\Res_s\bigl(\Phi_m^+(s),\widehat\Delta(w,s)\bigr).
\]
The order is part of the notation: the reversed order differs by
\((-1)^{3\deg\Phi_m^+}\) and is not the definition of \(R_m\).
If
\[
  \rho_m:=\frac{1}{\min\{\abs{w}:R_m(w)=0\}},
\]
then:
\begin{enumerate}[label=\textup{(\roman*)}]
  \item For every \(m\ge 4\),
  \[
    \deg_w R_m=3\varphi(2m).
  \]
  If \(m\) is even, then \(R_m(w)\) is an even polynomial.

  \item One has
  \[
    \rho_2=\sqrt{2+\sqrt3}.
  \]

  \item As \(m\to\infty\),
  \[
    \rho_m
    =3-\frac{11\pi^2}{17m^2}+O(m^{-4}).
  \]
  More precisely, the branch selection
  \(\rho_m=\rho(2\cos(\pi/m))\) is proved for every \(m\ge32\), and on
  that range \Cref{thm:rho-gap-m12} gives the expansion through order
  \(m^{-12}\). The small values displayed later are
  verified directly from the defining resultant rather than from this
  asymptotic branch-selection statement. In particular,
  \[
    3-\rho_m\sim \frac{11\pi^2}{17m^2}.
  \]
\end{enumerate}
\end{theorem}

The endpoint-selection criterion used to prove the large-\(m\) branch
selection is stated and proved in \Cref{thm:self-dual-endpoint-selection}.
For context, let \(B(u)=B_0+uB_1\) be a finite non-negative labelled
kernel for which \(B(1)\) is primitive, \(B(1)\mathbf1=q\mathbf1\), and
an involution exchanges \(B_0\) and \(B_1\).  If the associated completed
determinant \(H(w,s)\) satisfies the stated spectral-gap, endpoint
isolation, and effective root-separation hypotheses, then the real
cyclotomic root products
\[
  R_m(w):=\prod_{\alpha\in A_m}H(w,\alpha)
\]
select the endpoint slices \(\pm2\cos(\pi/m)\) for all sufficiently large
\(m\).  Any Taylor jet of the selected branch then gives the
corresponding large-\(m\) expansion.  In this paper this criterion is
used only through \Cref{cor:sync-endpoint-selection-application}, where
\Cref{lem:root-separation} verifies its hypotheses for the displayed
ten-state kernel with
\[
  \eta=\frac1{100},\qquad \varepsilon=\frac1{200},\qquad m_0=32.
\]

We also record the primitive-orbit asymptotics of the unweighted kernel.

\subsection{Previous work}

The rationality of dynamical \(\zeta\)-functions for shifts of finite
type is classical, beginning with Bowen--Lanford and Manning and
continuing through Parry--Pollicott and Bass
\cite{BowenLanford1970Zeta,Manning1971Axiom,ParryPollicott1983,Bass1992}.
We use this tradition only to place the determinant
\(\det(I-zB(u))\) in context; the determinant itself is computed
directly from the displayed matrix. General symbolic-dynamics
background is available in \cite{ArtinMazur1965,Ruelle1976,
LindMarcus2021,Kitchens1998}. The determinant manipulations also sit
near the multivariable graph-zeta determinant identities of
Foata--Zeilberger \cite{FoataZeilberger1999}. The standard
algebraic-geometry and function-field inputs are cited at the points
where they enter the proof: the smooth plane quartic genus formula,
quadratic ramification, and Riemann--Hurwitz.

The closest literature falls into three layers.

The word ``synchronisation'' is used here in the finite-automaton sense:
the motivating comparison construction has recurrent states organized by
how labelled paths re-align.  This places the example near the classical
literature on synchronizing automata and road colouring, including
\v{C}ern\'y's original work, Trahtman's solution of the road-colouring
problem, and Volkov's survey of the \v{C}ern\'y conjecture
\cite{Cerny1964,Trahtman2009,Volkov2008}.  The present paper does not
use any global synchronizing-automata theorem; the cited literature
explains the terminology and the finite-state motivation for isolating a
small recurrent kernel.

Classical background on graph zeta functions, coverings, and
Artin-style \(L\)-factorisations goes through Hashimoto
\cite{Hashimoto1990,Hashimoto1992}, Stark and Terras
\cite{StarkTerras1996,StarkTerras2000,StarkTerras2007}, and
Kotani--Sunada \cite{KotaniSunada2000}; see also Terras's monograph
\cite{Terras2011}. The voltage-graph and covering language for cyclic
lifts is classical \cite{GrossTucker2001}, while the twisted-operator
viewpoint relevant to the present Fourier decomposition connects to the
covering theory of Cimasoni--Kassel \cite{CimasoniKassel2023}. Spectral
properties of arbitrary lifts of graphs, directly relevant to
\Cref{prop:sync-cyclic-lift-factorisation}, are studied systematically
in \cite{DalfoFiolGarriga2019}.

Direct precedents for determinant expressions and covering
factorisations on digraphs and weighted graphs are due to Mizuno and
Sato \cite{MizunoSato2000,MizunoSato2001,MizunoSato2002,MizunoSato2004},
to Sato \cite{Sato2006}, and, on the dynamical/determinantal side, to
Morita and collaborators
\cite{Morita2020,IdeIshikawaMoritaSatoSegawa2021}. These works are the
closest antecedents to the cyclic-lift/resultant discussion in the
present note.

For the cyclotomic part, classical graph-zeta determinant formulas
explain why finite labelled kernels produce rational determinants;
graph-covering and Artin \(L\)-factorisation formulas explain why cyclic
lifts split into character factors; and real cyclotomic parametrisations
and cyclotomic-resultant conventions are standard
\cite{Apostol1970,CoxLittleOShea2005}.  The additional step needed in this case study is
effective endpoint selection: among all primitive real cyclotomic
slices, one must prove that no interior conjugate contributes a smaller
zero than the endpoint conjugate.  \Cref{thm:self-dual-endpoint-selection}
packages the elementary ingredients used here--self-dual completion,
strict unitary spectral gap, rational separation intervals, and the
endpoint Taylor jet--so that the ten-state verification can be audited
without repeating the selection argument at every \(m\).  The
substantive data remain the explicit certificates for this kernel.

Finite-group extensions of shifts of finite type provide the standard
language for tracking output labels in skew products. We use the
foundational framework of Adler--Kitchens--Marcus
\cite{AdlerKitchensMarcus1985} and Fiebig \cite{Fiebig1993}, together
with the \(\ZZ_q\)-extension viewpoint of Pollicott--Sharp
\cite{PollicottSharp1994}. The matrix-over-group-ring perspective of
Boyle--Schmieding \cite{BoyleSchmieding2017} is relevant here through
the cyclic Fourier factorisation of
\Cref{prop:sync-cyclic-lift-factorisation}. The underlying two-track
comparison can be viewed as a fibre product of labelled graphs, in the
sense of Nasu's textile systems \cite{Nasu1995}, but the parent
construction is not used in the proofs.  This article isolates the
ten-state recurrent kernel because its completed determinant remains
degree \(6\) while carrying nontrivial quotient geometry, generic
\(S_6\) Galois group, and explicit large-\(m\) cyclotomic asymptotics.

\subsection{Organisation}

\Cref{sec:preliminaries} recalls primitive orbit counts and defines the
weighted synchronisation kernel together with its self-duality.
\Cref{sec:completion-twists} contains the explicit determinant, the
cyclic-lift Fourier decomposition, the irreducibility and integrality
of the completed curve, its geometry, the resultant degree law, and the
asymptotics of the twisted Perron roots.
\Cref{sec:appendix-computation} records the explicit transition matrix,
the executable Sage certificate runner, and the exact
theorem-to-command map for the computer-assisted checks.







